\documentclass[10pt,a4paper]{amsart}
\usepackage[margin=19mm]{geometry}
\usepackage{amsmath,amssymb,mathtools}
\usepackage[scr=rsfs,cal=euler]{mathalfa}
\usepackage{xfrac}
\usepackage[unicode,hidelinks,bookmarks=false]{hyperref}
\newcommand{\ie}{i.e.\ }
\newcommand{\cf}{cf.\ }
\newcommand{\resp}{respectively\ }
\newcommand{\Subsection}{\subsection}
\providecommand{\nobreakdash}{\nobreak\hbox{-}}
\setlength{\emergencystretch}{2em}
\multlinegap0pt
\usepackage{bm}
\usepackage{bbm}
\usepackage{dsfont}
\usepackage{xspace}
\usepackage{cancel}
\usepackage{mathtools}
\usepackage{amsthm}
\makeatletter
\@ifundefined{theorem}{\newtheorem{theorem}{Theorem}}{}
\@ifundefined{proposition}{\newtheorem{proposition}[theorem]{Proposition}}{}
\makeatother
\title[Quasi-kernels in Hereditary Classes and Applications to Brea]{Quasi-kernels in Hereditary Classes and Applications to Break}
\author{Jiangdong Ai, Tianyu Huang, Xiangzhou Liu, Chaoliang Tang, Zongqun Xu}
\date{Source: arXiv:2606.16571v2 (CC BY 4.0)}
\begin{document}
\maketitle

\noindent\emph{Source and licence.} Quasi-kernels in Hereditary Classes and Applications to Break. Version arXiv:2606.16571v2 (CC BY 4.0), distributed under the Creative Commons Attribution 4.0 International licence, as recorded in the arXiv licence metadata for this exact version and shown on the abstract page https://arxiv.org/abs/2606.16571v2. Full rights notice: licenses/arxiv-2606.16571-CC-BY-4.0.txt.\par\smallskip

\noindent\textit{Editorial scope.} Counted unit: the Theorem whose source label is \texttt{thm:general-2n3} in the article (source0.tex lines 482--490, source label \texttt{thm:general-2n3}), delivered together with the article's own complete original proof of it (source0.tex lines 492--684, one \texttt{proof} environment of 193 lines). No mathematics is written, repaired or re-derived: statement and proof are the article's own text line for line. Required context retained uncounted: prop:sinks-general (lines 281--292). Every definition, display and label of those lines is retained unchanged. Omitted from this package and not redistributed: 1--280, 293--481, 491--491, 685--1058. Nothing inside a retained excerpt is omitted or altered: every retained body is delivered exactly as the article's lines are written, so no omission receipt of any kind accompanies this package. Citations: the retained excerpts carry no citation command at all, so no bibliography entry of the article is transcribed and no reference is redistributed. Reference adaptation: none was needed and none was made; every reference inside the delivered unit resolves inside it, so no unresolved reference or citation is produced. Printed slips kept literal: no printed word, symbol, hypothesis or number of the counted statement or of its proof was altered. Layout and class: the article's own class is \texttt{article} with packages including babel, geometry, amsmath, amssymb, amsthm, mathtools, xcolor, enumitem, hyperref, microtype; the wrapper is the lane's pinned \texttt{amsart} editorial wrapper at 10pt on A4, which restates the theorem-like environments the retained text needs and adds the article's own macro definitions that the retained text needs (none), with extra wrapper packages (bm, bbm, dsfont, xspace, cancel, mathtools). The delivered numbering is the unit's own. Assets: no retained span contains a figure environment, a graphics-inclusion command, a PDF-page inclusion, a table or a TikZ picture; no image file is copied into this package, the package declares no asset dependency and no image file is redistributed with it. Licence: version arXiv:2606.16571v2 of this article is distributed under the Creative Commons Attribution 4.0 International licence, as recorded in the arXiv licence metadata for this exact version and shown on the abstract page https://arxiv.org/abs/2606.16571v2. A full rights notice is delivered at licenses/arxiv-2606.16571-CC-BY-4.0.txt. Characters: every retained excerpt is pure ASCII, so the delivered engine, XeLaTeX with its default Latin Modern text font, reports no missing character.
\begin{proposition}\label{prop:sinks-general}
Let \(\mathcal G\) be a hereditary class satisfying the weighted
half-neighborhood property. Let \(D\in\mathcal G\), and let
\[
    Z=\{v\in V(D): v \text{ is a sink and not a source in }D\},
    \qquad z=|Z|.
\]
Then \(D\) has a quasi-kernel \(K\) such that
\[
    |K|\le |V(D)|-\left\lceil\frac z2\right\rceil .
\]
\end{proposition}

\begin{theorem}\label{thm:general-2n3}
Let \(\mathcal G\) be a hereditary class satisfying the weighted
half-neighborhood property.
Let \(D\in\mathcal G\) be source-free, and let \(n=|V(D)|\).
Then \(D\) has a quasi-kernel \(K\) such that
\[
    |K|\le \frac{2n}{3}.
\]
\end{theorem}

\begin{proof}
By Proposition~\ref{prop:sinks-general}, every induced subdigraph
\(H\in\mathcal G\) with \(z\) sinks that are not sources
has a quasi-kernel \(Q\) satisfying
\[
    |Q|
    \le
    |V(H)|-\left\lceil\frac z2\right\rceil .
\]

Choose an inclusion-minimal quasi-kernel \(K\subseteq V(D)\). Let
\[
    N=N_D^+(K),
    \qquad
    M=V(D)\setminus (K\cup N).
\]
Thus
\[
    V(D)=K\dot\cup N\dot\cup M.
\]

Consider the bipartite graph with parts \(K\) and \(N\), where
\(ku\) is an edge whenever \(k\to u\) in \(D\).
Let \(\mathcal M\) be a maximum matching.
Let \(K_1\subseteq K\) be the set of vertices saturated by \(\mathcal M\),
and let
\[
    K_2=K\setminus K_1.
\]

Write
\[
    r=|K_1|,
    \qquad
    s=|K_2|,
    \qquad
    p=|N|,
    \qquad
    m=|M|.
\]
Then
\[
    n=r+s+p+m,
\]
and since \(\mathcal M\) is a matching,
\[
    r\le p.
\]

\medskip
\noindent\textbf{Claim 1.}
Every vertex \(u\in N\) has an in-neighbor in \(K_1\).

\begin{proof}
By definition of \(N=N_D^+(K)\), every \(u\in N\) has an in-neighbor in \(K\).
If none of these in-neighbors belonged to \(K_1\), then one of them would lie
in \(K_2\), and matching this vertex to \(u\) would enlarge \(\mathcal M\),
a contradiction.
\end{proof}

\medskip
\noindent\textbf{Claim 2.}
No vertex of \(K_2\) has an in-neighbor in \(N\).

\begin{proof}
Suppose, to the contrary, that \(u\to v\) for some \(u\in N\) and
\(v\in K_2\). By Claim~1, there exists \(w\in K_1\) such that \(w\to u\).
Then
\[
    w\to u\to v,
\]
so \(v\) is reached from \(K\setminus\{v\}\) within distance at most 2.

We claim that \(K\setminus\{v\}\) is still a quasi-kernel of \(D\),
contradicting the inclusion-minimality of \(K\).
The set \(K\setminus\{v\}\) remains independent.
Moreover, every vertex of \(N_D^+(v)\) lies in \(N\), and every vertex of \(N\)
has an in-neighbor in \(K_1\) by Claim~1. Hence all former out-neighbors of
\(v\) are still reached from \(K\setminus\{v\}\) within distance at most \(1\),
while \(v\) itself is reached by the directed path \(w\to u\to v\).
Thus \(K\setminus\{v\}\) is a quasi-kernel of \(D\), a contradiction.
\end{proof}

Now consider
\[
    H=D[K_2\cup M].
\]
Since \(\mathcal G\) is hereditary and \(D\in\mathcal G\),
we have \(H\in\mathcal G\).

Every vertex of \(K_2\) is a sink in \(H\), because \(K\) is independent and
all out-neighbors of vertices of \(K\) lie in \(N\).
Moreover, no vertex of \(K_2\) is a source in \(H\).
Indeed, since \(D\) is source-free, every vertex of \(K_2\) has an
in-neighbor in \(D\). This in-neighbor cannot lie in \(K\), by the
independence of \(K\), and cannot lie in \(N\), by Claim~2.
Hence it lies in \(M\).

Therefore \(H\) has at least \(s=|K_2|\) vertices that are sinks but not
sources. By Proposition~\ref{prop:sinks-general}, \(H\) has a quasi-kernel
\(K'\) such that
\[
    |K'|
    \le
    |V(H)|-\left\lceil\frac s2\right\rceil
    \le
    m+\frac s2.
\]

Define
\[
    K^*
    =
    K'
    \cup
    \bigl(K_1\setminus N_D^+(K')\bigr).
\]

We claim that \(K^*\) is a quasi-kernel of \(D\).

First, \(K^*\) is independent.
The set \(K'\) is independent, and
\(K_1\setminus N_D^+(K')\) is independent because it is contained in \(K\).
There is no arc from \(K'\) to
\(K_1\setminus N_D^+(K')\) by definition.
There is also no arc from \(K_1\setminus N_D^+(K')\) to \(K'\).
Indeed, if \(x\in K_1\) sent an arc to some \(y\in K'\), then
\(y\notin K_2\), since \(K\) is independent, and
\(y\notin M\), since then \(y\in N_D^+(K)\), contradicting \(y\in M\).

Next, \(K^*\) 2-dominates \(D\).
Since \(K'\) is a quasi-kernel of \(H=D[K_2\cup M]\),
every vertex of \(K_2\cup M\) is reached from \(K'\) within distance at most 2.

Every vertex of \(K_1\) is reached from \(K^*\) within distance at most \(1\):
either it belongs to \(K_1\setminus N_D^+(K')\), or it lies in \(N_D^+(K')\).

Finally, if \(u\in N\), then by Claim~1 there exists \(w\in K_1\)
with \(w\to u\).
Since \(w\) is reached from \(K^*\) within distance at most \(1\),
the vertex \(u\) is reached from \(K^*\) within distance at most 2.
Thus \(K^*\) is a quasi-kernel of \(D\).

Moreover,
\[
    |K^*|
    \le
    |K'|+|K_1|
    \le
    m+\frac s2+r.
\]

It remains to compare \(K\) and \(K^*\).
Suppose, for a contradiction, that both have size greater than \(2n/3\).

Since \(|K|=r+s\), we have
\[
    r+s>\frac{2(r+s+p+m)}{3}.
\]
Thus
\[
    r+s>2p+2m.
\]
Using \(p\ge r\), this gives
\[
    s>2m.
\]

On the other hand, since
\[
    |K^*|\le r+m+\frac s2
\]
and \(|K^*|>2n/3\), we obtain
\[
    r+m+\frac s2>\frac{2(r+s+p+m)}{3}.
\]
Equivalently,
\[
    2r+2m-s>4p.
\]
Using \(p\ge r\), this implies
\[
    2r+2m-s>4r,
\]
and hence
\[
    s<2m-2r\le 2m.
\]
This contradicts \(s>2m\).

Therefore at least one of \(K\) and \(K^*\) has size at most \(2n/3\).
Since both are quasi-kernels of \(D\), the result follows.
\end{proof}

\end{document}
